% This is text associated with the open textbook "Open Electromagnetics"
% See immediately below for author, date
% This work is licensed under the Creative Commons Attribution-ShareAlike 4.0 International License. (CC BY SA 4.0) 
% To view a copy of the license, visit http://creativecommons.org/licenses/by-sa/4.0/

%ORB TITLE Skin Depth
%ORB AUTHOR 0        # S.W. Ellingson, Virginia Tech (ellingson@vt.edu)
%ORB DATE 20191116   # yyyymmdd
%ORB ABSTRACT

\label{m0158_Skin_Depth} % <-- Must match name of this file and module directory name.  Don't mess with this!
\index{skin depth|(}

The electric and magnetic fields of a wave are diminished as the wave propagates through lossy media.
The magnitude of these fields is proportional to
%
\begin{equation}
e^{-\alpha l} \nonumber
\end{equation}
%
where $\alpha\triangleq\mbox{Re}\left\{\gamma\right\}$ is the attenuation constant
\index{propagation constant!attenuation}
(SI base units of m$^{-1}$), $\gamma$ is the propagation constant, and $l$ is the distance traveled.
Although the rate at which magnitude is reduced is completely described by $\alpha$, particular values of $\alpha$ typically do not necessarily provide an intuitive sense of this rate.

An alternative way to characterize attenuation is in terms of \emph{skin depth}
\index{skin depth}
$\delta_s$, which is defined as the distance at which the magnitude of the electric and magnetic fields is reduced by a factor of $1/e$.  
In other words:
%
\begin{equation}
e^{-\alpha\delta_s} = e^{-1} \cong 0.368
\label{m0158_esddef}
\end{equation}
%
%%%%%%%%%%%%%%%%%%%%%%%%%%%%%%%%%%%%%%%%%%%%%%%
%ORB HIGHLIGHT
\begin{mdframed}[backgroundcolor=yellow!20,nobreak=true] 
Skin depth $\delta_s$ is the distance over which the magnitude of the electric or magnetic field is reduced by a factor of $1/e \cong 0.368$.
\end{mdframed}
%%%%%%%%%%%%%%%%%%%%%%%%%%%%%%%%%%%%%%%%%%%%%%%
%
Since power is proportional to the square of the field magnitude, $\delta_s$ may also be interpreted as the distance at which the power in the wave is reduced by a factor of $(1/e)^2\cong 0.135$.  
In yet other words: $\delta_s$ is the distance at which $\cong 86.5$\% of the power in the wave is lost.

This definition for skin depth makes $\delta_s$ easy to compute:  From Equation~\ref{m0158_esddef}, it is simply
%
\begin{equation}
\boxed{
\delta_s = \frac{1}{\alpha}
}
\end{equation}

The concept of skin depth is most commonly applied to good conductors.
For a good conductor, $\alpha\approx \sqrt{\pi f\mu\sigma}$ (Section~\ref{m0157_Good_Conductors}), so 
%
\begin{equation}
\delta_s \approx \frac{1}{\sqrt{\pi f\mu\sigma}} ~~~\mbox{(good conductors)}
\label{m0158_eSDGC}
\end{equation}

%%%%%%%%%%%%%%%%%%%%%%%%%%%%%%%%%%%%%%%%%%%%%%%
%ORB EXAMPLE
~\\ \begin{example} 
\begin{mdframed}[backgroundcolor=brown!20,linecolor=brown!20]\raggedright
Skin depth of aluminum. 
\index{aluminum}
\\~\\
Aluminum, a good conductor, exhibits $\sigma\approx 3.7 \times 10^7$~S/m and $\mu\approx\mu_0$ over a broad range of radio frequencies.
Using Equation~\ref{m0158_eSDGC}, we find $\delta_s \sim 26~\mu$m at 10~MHz.
Aluminum sheet which is 1/16-in ($\cong 1.59$~mm) thick can also be said to have a thickness of $\sim 61\delta_s$ at 10~MHz.
The reduction in the power density of an electromagnetic wave after traveling through this sheet will be
$$ \sim\left(e^{-\alpha\left(61\delta_s\right)}\right)^2 = \left(e^{-\alpha\left(61/\alpha\right)}\right)^2 = e^{-122} $$
which is effectively zero from a practical engineering perspective.  Therefore, 1/16-in aluminum sheet provides excellent shielding from electromagnetic waves at 10~MHz.

\noindent
At 1~kHz, the situation is significantly different.
At this frequency, $\delta_s\sim 2.6$~mm, so 1/16-in aluminum is only $\sim 0.6\delta_s$ thick.
In this case the power density is reduced by only 
$$ \sim\left(e^{-\alpha\left(0.6\delta_s\right)}\right)^2 = \left(e^{-\alpha\left(0.6/\alpha\right)}\right)^2 = e^{-1.2} \approx 0.3 $$
This is a reduction of only $\sim70$\% in power density.
Therefore, 1/16-in aluminum sheet provides very little shielding at 1~kHz.
\end{mdframed}
% note figures will need to go between "\end{mdframed}" and "\end{example}"
\end{example}
%%%%%%%%%%%%%%%%%%%%%%%%%%%%%%%%%%%%%%%%%%%%%%%

%%%%%%%%%%%%%%%%%%%%%

%{\bf Additional Reading:}
%\begin{itemize}
%\item ``\href{https://en.wikipedia.org/wiki/Skin_effect}{Skin effect}'' on Wikipedia.
%\end{itemize}

\index{skin depth|)}


